\documentclass[11pt]{article}
\usepackage{fontspec}
\setmainfont{Times New Roman}
\setsansfont{Arial}
\setmonofont{Consolas}
\usepackage{amsmath,amssymb,mathtools}
\usepackage{geometry}
\usepackage{microtype}
\geometry{a4paper,margin=2.28cm}
\setlength{\parindent}{1.15em}
\setlength{\parskip}{0pt}
\providecommand{\mo}{\mathfrak{o}}
\providecommand{\mc}{\mathfrak{c}}
\providecommand{\mM}{\mathfrak{M}}
\emergencystretch=140em
\allowdisplaybreaks
\sloppy
\hbadness=10000
\vbadness=10000
\hfuzz=2pt
\vfuzz=10pt
\raggedbottom
\begin{document}
% Унит: Папер34 Сецтион24 соурце-фиделиты цонтинуатион в001.
% Лангуаге лане: Интерславиц Латин.
% Дирецт соурце: соурцес/папер34/соурце_фиделиты/Noether_Папер34_Сецтион24_ОРИГИНАЛ_СЦАН_ВИТНЕСС_в001.текс.
\section*{Emmy Noether: Хиперкомплексне величины и теорија представјень}
\emph{Праца 34, §24. Изворно-верны превод из сканованого немецкого сведка.}

\subsection*{§ 24. Следы и карактере}

Ако в представјеньу $\mo\sim\mathfrak D$ хиперкомплексного система $\mo$ елементу $a$ јест придана матрица $A$, тогда ставимо
\[
        \operatorname{Sp}_{D}(a)=\operatorname{Sp} A.
\]
Следы сут линеарне функције:
\[
        \operatorname{Sp}(c+d)=\operatorname{Sp}(c)+\operatorname{Sp}(d),
        \qquad
        \operatorname{Sp}(c\alpha)=\alpha\operatorname{Sp}(c).
\]
Еквивалентне представјеньа имајут те саме следы.

След в редуцибилном представјеньу јест сума следов в представјеньах, кторе сут дане композицијными факторами.

Доказ. При подобном избору базы матрица има блокову форму
\[
        \begin{pmatrix}
        A_{11}&0&\cdots&0\\
        *&A_{22}&\cdots&0\\
        \vdots&\vdots&\ddots&\vdots\\
        *&*&\cdots&A_{ss}
        \end{pmatrix},
\]
зато јеј след јест
\[
        \operatorname{Sp}(A_{11})+\operatorname{Sp}(A_{22})+\cdots+\operatorname{Sp}(A_{ss}).
\]

Главны след именује се след при регуларном представјеньу. Редукованы след именује се сума следов при различнох иредуцибилных представјеньах.

Ако $c$ јест елемент максималного нилпотентного идеала $\mc$, тогда $\operatorname{Sp}(c)=0$ за всако представјенье.

Доказ. Достатно јест изучати представјеньа чрез просте леве идеалы $\mo/\mc$, бо в всаком другом представјеньу појавјајут се само ти композицијне факторы, а след складаје се адитивно. Але $c$ анулује все елементы $\mo/\mc$; зато всакому $c$ из $\mc$ принадлежи нулова матрица, и зато $\operatorname{Sp}(c)=0$.

Ако $\mo$ јест двустранно просты и $P$ алгебраично затворјено, значи матрично колцо
\[
        \mo=\sum c_{ik}P,
\]
тогда в иредуцибилном представјеньу $\mo$ елементу
\[
        a=\sum c_{ik}\alpha_{ik}
\]
принадлежи матрица $(\alpha_{ik})$. Зато
\[
        \operatorname{Sp}_{\rm red}(a)=\sum_i\alpha_{ii},
        \qquad
        \operatorname{Sp}_{\rm pr}(a)=n\,\operatorname{Sp}_{\rm red}(a)=n\sum_i\alpha_{ii}.
\]
Особливо
\[
        \operatorname{Sp}_{\rm red}(c_{ik})=0\quad(i\ne k),\qquad
        \operatorname{Sp}_{\rm red}(c_{ii})=e,
\]
и такоже
\[
        \operatorname{Sp}_{\rm pr}(c_{ik})=0\quad(i\ne k),\qquad
        \operatorname{Sp}_{\rm pr}(c_{ii})=n e.
\]

При преходу од $P$ к алгебраично затворјеному польу $\Omega$ главны след не мјеньа се, бо може се израчунати из тој самеј базы. То не јест правда за редукованы след.

Следы елементов $a$ система без радикала в абсолутно иредуцибилных представјеньах именујут се карактере и означајут се $\chi(a)$, или, ако треба указати дано представјенье, $\chi^{(\nu)}(a)$.

В иредуцибилном представјеньу степена $n_\nu$ елементы центра, по §21, представјајут се диагоналными матрицами $E\theta^{(\nu)}$, кде $z\mapsto\theta^{(\nu)}(z)$ јест представјенье првого степена центра, или хомоморфизм центра в $\Omega$. Зато след има вредност $n_\nu\theta^{(\nu)}(z)$. Так хомоморфизмы $\theta$ центра сут связане с карактерами $\chi$ релацијеју
\[
        \chi^{(\nu)}(z)=n_\nu\theta^{(\nu)}(z).                \tag{1}
\]

В комутативном случају $n_\nu=1$, и карактере саме давајут хомоморфизмы (ср. §22).

Ако полье $\Omega$ има карактеристику нула, что будемо в дальньем всагда предполагати, тогда можно (1) разделити прзез $n_\nu$:
\[
        \theta^{(\nu)}(z)=\frac{\chi^{(\nu)}(z)}{n_\nu}.
\]
Хомоморфна властност за $\theta^{(\nu)}$ изражаје се формулами
\[
        \frac{\chi^{(\nu)}(z+z')}{n_\nu}
        =
        \frac{\chi^{(\nu)}(z)}{n_\nu}
        +
        \frac{\chi^{(\nu)}(z')}{n_\nu},
\]
\[
        \frac{\chi^{(\nu)}(zz')}{n_\nu}
        =
        \frac{\chi^{(\nu)}(z)}{n_\nu}
        \frac{\chi^{(\nu)}(z')}{n_\nu}.
\]

Класа представјень јест уже једнозначно определьена самыми следами матриц. Да бы знати следы всех матриц, јасно достатно јест знати следы базовых елементов в даном представјеньу.

Доказ. Модул представјеньа $R$ јест познаны, ако се зна, колико разов всакы иредуцибилны модул представјеньа $\mM_\nu$ појавја се в ньем како директны суманд. Ако та число јест $\mu_\nu$, тогда след $e^{(\nu)}$ в представјеньу јест $\mu_\nu n_\nu$, кде $n_\nu$ јест степень иредуцибилного представјеньа. Зато
\[
        \mu_\nu=\frac{\operatorname{Sp}(e^{(\nu)})}{n_\nu}.
\]

\end{document}

